\section{Chương~\ref{sec:livingmachine}. Cỗ máy từ nơ-ron sống}

Hành vi của mẫu nuôi cấy trên mảng điện cực đã được đo trong các thí nghiệm trực tiếp có ghi và kích thích; cơ chế của các đợt bùng phát dựa trên mô hình cạn kiệt khớp thần kinh và thí nghiệm trên vi mẫu nuôi cấy, còn các khẳng định về dopamine trong đĩa nuôi dựa trên những thí nghiệm đơn lẻ, mà một phần trong số đó chính tác giả cũng chỉ coi là minh họa.

{\footnotesize
\setlength{\tabcolsep}{3pt}\renewcommand{\arraystretch}{1.15}
\begin{longtable}{>{\raggedright}p{0.5cm}>{\raggedright}p{4.0cm}>{\raggedright}p{5.6cm}>{\raggedright}p{2.3cm}>{\raggedright\arraybackslash}p{1.7cm}}
\toprule
TT & Khẳng định & Thí nghiệm và điều đã đo & Nguồn & Trạng thái \\
\midrule\endfirsthead
\toprule
TT & Khẳng định & Thí nghiệm và điều đã đo & Nguồn & Trạng thái \\
\midrule\endhead
1 & Mẫu nuôi cấy trên chip: phần lớn là \nt{GLUT}, một tỷ lệ nhỏ hơn \nt{GABA} tùy mẫu; trong mẫu nuôi cấy hồi hải mã khoảng $6\%$ & Mạng hàng nghìn nơ-ron trên mảng điện cực là thí nghiệm tiêu chuẩn (dòng~3). Tỷ lệ nơ-ron \nt{GABA} trong mẫu nuôi cấy hồi hải mã được đếm bằng nhuộm GABA: khoảng $6\%$ số nơ-ron. Với mẫu nuôi cấy vỏ não, chúng tôi không đưa ra con số đã được kiểm chứng & \cite{Benson1994} & đã đo \\
2 & Organoid: mô thần kinh ba chiều cỡ khoảng một milimét từ tế bào gốc người & Từ tế bào gốc đã nuôi được organoid ba chiều có nhiều vùng não, trong đó có vỏ não với các lớp tế bào tiền thân và nơ-ron trưởng thành & \cite{Lancaster2013} & đã đo \\
3 & Không có đầu vào, mẫu nuôi cấy bùng lên thành từng đợt, cách nhau từ một giây đến vài phút tùy mẫu & $58$ mẫu nuôi cấy vỏ não mật độ từ $3\,000$ đến $50\,000$ nơ-ron được ghi trong năm tuần ($963$ bản ghi nửa giờ): sau hai tuần, ở gần như mọi mẫu, các đợt bùng phát toàn mạng chiếm ưu thế; khoảng cách giữa các đợt từ $1$ đến $300$~s và phụ thuộc mạnh vào mẫu & \cite{Wagenaar2006} & đã đo \\
4 & Đợt bùng phát kết thúc do cạn chất dẫn truyền, khoảng cách $T_{\text{đợt}}\approx\tau_{\text{ph}}\ln\frac{1}{1-x^*}$~\eqref{eq:burstinterval}; $2$--$4$~s là ước tính mô hình với $\tau_{\text{ph}}=0.8$~s, sự hồi phục đo được thì chậm hơn & Công thức được suy ra trong chương. Mô hình: mạng có khớp thần kinh bị cạn kiệt tự tạo ra các đợt bùng phát toàn mạng. Thí nghiệm trên vi mẫu nuôi cấy $4$--$30$ nơ-ron: độ dài đợt bùng phát phụ thuộc vào thời gian từ đợt trước, làm cạn túi chất dẫn truyền thì xóa các đợt; kết luận của tác giả là đợt bùng phát bị giới hạn bởi kho chất dẫn truyền. Kho hồi phục ở đó mất vài chục giây, và với cùng công thức thì cho khoảng cách vài chục giây đến vài phút & suy ra trong chương; \cite{Tsodyks2000,CohenSegal2011,Tsodyks1997} & lý thuyết \\
5 & ``Người thầy im lặng'': kích thích bị tắt khi đáp ứng mong muốn xuất hiện, và đáp ứng được củng cố & Mẫu nuôi cấy được kích thích ở tần số $0.3$--$1$~Hz cho đến khi đáp ứng đã chọn xuất hiện sau kích thích $50\pm10$~ms, rồi tắt kích thích trong $5$~phút; sau vài chu kỳ, đáp ứng mong muốn được gây ra ngay; đã vẽ các đường cong học & \cite{Shahaf2001} & đã đo \\
6 & Mẫu nuôi cấy chơi quần vợt; mức tăng có ý nghĩa của các chuỗi trong một phiên chỉ có khi phản hồi đầy đủ & Chi tiết ở chương~\ref{sec:dishbrain} và phụ lục của nó: với im lặng thay cho kích thích, cuối phiên mẫu nuôi cấy cũng chơi tốt hơn so với không có phản hồi, nhưng hiệu ứng nhỏ hơn; việc học từ phiên này sang phiên khác không bền vững & \cite{Kagan2022} & đã đo \\
7 & Việc mẫu nuôi cấy học để dự đoán đầu vào của mình là giả thuyết; những lời lớn về sự ``có tri giác'' bị phản bác & Nguyên lý năng lượng tự do là một đề xuất lý thuyết; chưa có thí nghiệm trực tiếp nào phân biệt nó với các cách giải thích đơn giản hơn. Ba mươi nhà nghiên cứu trong một bài phản hồi chỉ ra rằng sự thay đổi hành vi trong vòng lặp không chứng tỏ trí tuệ hay cảm giác & \cite{Friston2010,Balci2023} & giả thuyết \\
8 & Mẫu nuôi cấy điều khiển cơ thể ảo tới mục tiêu khi chọn đúng mẫu kích thích & Chương trình tự chọn các mẫu kích thích huấn luyện theo kết quả hiện tại; trong vài chục phút, mạng đạt được các trạng thái mong muốn, tính dẻo duy trì $80$~phút sau $2$~giờ huấn luyện. Cùng những kích thích ấy, đưa vào không gắn với hành vi, thì không có tác dụng & \cite{Bakkum2008} & đã đo \\
9 & Hồ chứa: chỉ lớp đọc ra tuyến tính $y=W_{\text{ra}}\mathbf r$ được huấn luyện~\eqref{eq:reservoir} & Lý thuyết tính toán hồ chứa: một hệ phi tuyến bất kỳ có trí nhớ tắt dần cộng với một đầu ra tuyến tính được huấn luyện & \cite{Maass2002,JaegerHaas2004} & lý thuyết \\
10 & Organoid làm hồ chứa phân biệt người nói với độ chính xác khoảng $78\%$ (theo tác giả); lớp đọc ra học theo sai số, còn liên kết của organoid thay đổi không cần thầy & Âm thanh tiếng nói của vài người được đưa vào organoid dưới dạng các mẫu dòng điện trên mảng điện cực; lớp đọc ra đã huấn luyện phân biệt người nói với độ chính xác khoảng $78\%$, theo báo cáo của tác giả. Họ cũng báo cáo rằng trong lúc huấn luyện, độ liên kết chức năng của organoid thay đổi, và gọi đó là học không giám sát bên trong organoid & \cite{Cai2023} & đã đo \\
11 & Một triệu nơ-ron tốn cho xung khoảng $10^{-4}$~W~\eqref{eq:dishpower} & Ước tính: giá một xung $10^{-10}$~J từ ngân sách năng lượng của vỏ não~\eqref{eq:energy}, nhân với số xung; chưa có phép đo trực tiếp công suất của mẫu nuôi cấy & suy ra trong chương; \cite{Attwell2001} & lý thuyết \\
12 & Dopamine theo chớp sáng: $1$~mM dopamine trong lồng, $240$ lần lặp, số xung được gây ra tăng lên & Trên nền tảng organoid từ xa, dopamine trong lồng nhạy sáng ($1$~mM) được giải phóng bằng tia cực tím ($365$~nm, $0.8$~s) khi kích thích gây ra nhiều xung hơn; qua $240$ bước (khoảng $5$~giờ), số xung nhìn chung tăng. Tác giả viết rằng chỉ từ một thí nghiệm thì không thể xác lập mối liên hệ với dopamine; không có đối chứng & \cite{Jordan2024} & giả thuyết \\
13 & Dopamine từ tế bào sống thay đổi trọng số của transistor hữu cơ lâu dài và thuận nghịch & Tế bào tiết dopamine được nuôi phía trên một phần tử neuromorphic hữu cơ; dopamine bị oxy hóa ở điện cực làm thay đổi độ dẫn của kênh; đã cho thấy thay đổi trọng số lâu dài và sự đưa về lại & \cite{Keene2020} & đã đo \\
14 & Có những organoid với nơ-ron \nt{DOPA} hoạt động; dopamine của chúng chưa được dùng làm thầy dạy & Trong organoid từ tế bào gốc người đã tìm thấy nơ-ron \nt{DOPA} trưởng thành hoạt động điện và sự sản sinh dopamine. Chúng tôi không tìm thấy trong tài liệu thí nghiệm nào dùng dopamine này dạy một mạng khác & \cite{Jo2016} & đã đo \\
15 & Organoid giữ thăng bằng cây sào: kích thích dạy do chương trình chọn tốt hơn ngẫu nhiên, không được củng cố, qua khớp thần kinh \nt{GLUT} & Organoid vỏ não chuột trong vòng kín với bài toán ``cây sào trên xe đẩy'': với phần lớn organoid, tín hiệu do học tăng cường chọn cho kết quả tốt hơn tín hiệu ngẫu nhiên hoặc không có tín hiệu; sau $45$~phút nghỉ, sự cải thiện không còn; chặn thụ thể AMPA và NMDA làm mất nó & \cite{Robbins2026} & đã đo \\
16 & Không có thanh ghi điều khiển, mạng trên bàn thử không tự quyết định được thế nào là thành công (mệnh đề~\ref{prop:livingmachine}) & Trong mọi thí nghiệm ở các dòng 5--15, tín hiệu dạy do người làm thí nghiệm hoặc chương trình đặt ra; chưa có thí nghiệm nào mà mẫu nuôi cấy tự hình thành tiêu chí thành công --- đây là kết luận từ sự vắng mặt, không phải thí nghiệm & --- & giả thuyết \\
\bottomrule
\end{longtable}}
